% discussion/part_intro-DE.tex
\chapter*{Diskussion}
\phantomsection
\addcontentsline{toc}{chapter}{Diskussion}

Dieser Teil legt dar, wie das vollendete AIM-Framework in der weiteren
Ingenieurlandschaft einzuordnen ist.

\section*{Was wissen wir bereits?}

Die vorangehenden Teile haben AIM als vollständiges Alignment-basiertes
Informationsframework entwickelt.

Das Framework umfasst nun Geometrie, Zustände, Laufzeitauswertung,
Realität, Modellierung, Optimierung, Ingenieurbedingungen und praktische
Anwendungen.

Gemeinsam bilden diese Komponenten eine kohärente theoretische und
praktische Grundlage für die Modellierung von Eisenbahn-Alignments.

Das Framework selbst ist damit weitgehend vollständig.

\section*{Was fehlt noch?}

Auch ein vollständiges Framework bedarf der Interpretation.

Neue Begriffe müssen zu bestehenden Ingenieurtraditionen in Beziehung
gesetzt werden. Historische Ansätze sind neu zu bewerten. Zugrunde
liegende Annahmen sind zu untersuchen. Konsequenzen jenseits des
unmittelbaren Umfangs des Frameworks sind zu bedenken.

Das noch fehlende Element ist somit die Reflexion.

Die Frage lautet nicht länger, wie AIM funktioniert, sondern was AIM im
weiteren Kontext von Eisenbahnwesen und Informationsmodellierung bedeutet.

\section*{Was folgt?}

Dieser Teil eröffnet eine interpretierende Perspektive auf das
AIM-Framework. Der Schwerpunkt verschiebt sich von dessen Konstruktion zu
seinem Verständnis.

Die folgenden Kapitel setzen AIM in Beziehung zu:
\begin{itemize}
\item historischen Ansätzen des Eisenbahnwesens,
\item etablierten Modellierungstraditionen,
\item früheren Optimierungssystemen
\item und weiter reichenden Folgerungen für die Alignment-basierte
Informationsmodellierung.
\end{itemize}

Besondere Aufmerksamkeit gilt der Neuinterpretation früherer Begriffe aus
der Perspektive, die diese Arbeit entwickelt hat. Das Ziel ist nicht nur
der Vergleich, sondern das Verständnis.

\section*{Bezug zur AIM-Roadmap}

Innerhalb der AIM-Roadmap entspricht dieser Teil dem Übergang von der
Konstruktion zur Interpretation des Frameworks.

Die Roadmap-Skizze behandelt die zentrale Frage: Wie gelangt die Arbeit
von der technischen Vollendung zur begrifflichen Interpretation?

\begin{center}
\begin{tikzpicture}[
  node distance=1.4cm,
  box/.style={draw, rounded corners, minimum width=5.0cm, minimum height=0.9cm, align=center}
]
\node[box] (framework) {Vollständiges Framework};
\node[box, below=of framework] (reflection) {Reflexion};
\node[box, below=of reflection] (interpretation) {Interpretation};
\draw[->, thick] (framework) -- (reflection);
\draw[->, thick] (reflection) -- (interpretation);
\end{tikzpicture}
\end{center}

Die Entwicklung
\[
\text{vollständiges Framework}
\rightarrow
\text{Reflexion}
\rightarrow
\text{Interpretation}
\]
kennzeichnet den Übergang von der Entwicklung von AIM zum Verständnis
seiner Stellung in der weiteren Landschaft des Eisenbahnwesens und der
Informationsmodellierung.

Die Diskussion dient somit als begrifflicher Spiegel des Frameworks und
legt seine Annahmen, Folgerungen und Beziehungen zur etablierten Praxis
offen.

\begin{principle}[Interpretationsprinzip]
\label{princ:interpretation-principle}
Ingenieurframeworks sollen in AIM nicht nur konstruiert und angewandt,
sondern auch im Verhältnis zu historischer Praxis, alternativen Ansätzen
und ihren weiter reichenden begrifflichen Konsequenzen interpretiert
werden.
\end{principle}

\begin{summary}
Anwendungen zeigen den Nutzen.

Die Diskussion schafft Verständnis.

Dieser Teil reflektiert daher das AIM-Framework, seine Ursprünge, seine
Konsequenzen und seinen Platz in der weiteren Entwicklung des
Eisenbahnwesens und der Informationsmodellierung.
\end{summary}

Diese interpretierende Konsolidierung bereitet den Ausblick vor, in dem
dasselbe Framework auf künftige Forschungs- und Einsatzpfade projiziert
wird.
